This survey introduces reinforcement learning methods to researchers in the social sciences. The curse of dimensionality limits exact dynamic programming, forcing reliance on tractable reductions that a growing class of economic models resists. Reinforcement learning algorithms extend dynamic programming to problems with high-dimensional states, continuous actions, and strategic interactions through sample-based approximation. I review the theory connecting classical planning to modern learning algorithms, covering value-based methods, policy gradients, and their convergence properties. Applications include structural estimation of dynamic models, equilibrium computation in games of imperfect information, and preference learning from human feedback. I also examine the practical vulnerabilities of these methods, including brittleness, sample inefficiency, sensitivity to hyperparameters, and the absence of convergence guarantees outside of tabular settings. A companion survey \citep{RustRawat2026} covers the inverse problem of inferring preferences from observed behavior. Simulation code is publicly available.\footnote{\url{https://github.com/rawatpranjal/survey-of-reinforcement-learning-in-economics}}
